\documentclass[12pt]{article}
\usepackage{amsmath,amssymb,amsfonts}
\usepackage[margin=1in]{geometry}
\usepackage{bm}

\newcommand{\vect}[1]{\mathbf{#1}}

\begin{document}

\section*{Problem 3}

\textbf{Problem:} Prove from simple vector arguments and the conservation laws of mechanics that in an elastic collision between two particles of equal mass, one initially at rest, the velocity vectors after collision are perpendicular (if nonzero).

\bigskip
\textbf{Solution:}

Let both particles have mass $m$. Before the collision:
\begin{itemize}
\item Particle 1 has velocity $\vect{v}_0$.
\item Particle 2 is at rest.
\end{itemize}

After the collision:
\begin{itemize}
\item Particle 1 has velocity $\vect{v}_1$.
\item Particle 2 has velocity $\vect{v}_2$.
\end{itemize}

\textbf{Conservation of momentum:}
\begin{equation}
m\vect{v}_0 = m\vect{v}_1 + m\vect{v}_2 \quad \Longrightarrow \quad \vect{v}_0 = \vect{v}_1 + \vect{v}_2. \label{eq:mom}
\end{equation}

\textbf{Conservation of kinetic energy} (elastic collision):
\begin{equation}
\tfrac{1}{2}m|\vect{v}_0|^2 = \tfrac{1}{2}m|\vect{v}_1|^2 + \tfrac{1}{2}m|\vect{v}_2|^2 \quad \Longrightarrow \quad |\vect{v}_0|^2 = |\vect{v}_1|^2 + |\vect{v}_2|^2. \label{eq:energy}
\end{equation}

Now, from Eq.~\eqref{eq:mom}, square both sides using the dot product:
\begin{align}
|\vect{v}_0|^2 &= (\vect{v}_1 + \vect{v}_2) \cdot (\vect{v}_1 + \vect{v}_2) \\
&= |\vect{v}_1|^2 + 2\,\vect{v}_1 \cdot \vect{v}_2 + |\vect{v}_2|^2. \label{eq:squared}
\end{align}

Substituting Eq.~\eqref{eq:energy} into the left side:
\[
|\vect{v}_1|^2 + |\vect{v}_2|^2 = |\vect{v}_1|^2 + 2\,\vect{v}_1 \cdot \vect{v}_2 + |\vect{v}_2|^2.
\]

This immediately gives:
\[
2\,\vect{v}_1 \cdot \vect{v}_2 = 0 \quad \Longrightarrow \quad \boxed{\vect{v}_1 \cdot \vect{v}_2 = 0}.
\]

Therefore, if both $\vect{v}_1$ and $\vect{v}_2$ are nonzero, the velocity vectors after collision are perpendicular. \hfill $\blacksquare$

\textit{Note:} The degenerate cases are: (i) a head-on collision where particle 1 stops ($\vect{v}_1 = 0$, $\vect{v}_2 = \vect{v}_0$), or (ii) a miss ($\vect{v}_1 = \vect{v}_0$, $\vect{v}_2 = 0$). In both cases one velocity is zero and the perpendicularity condition is trivially satisfied (or undefined).

\end{document}
